\documentclass[a4paper,12pt,english]{article}
\usepackage{glossaries}
\usepackage[T1]{fontenc}
\usepackage{babel}
\usepackage{graphicx}
\usepackage[table,xcdraw]{xcolor}
\usepackage{hyperref}
\usepackage{blindtext}
\usepackage{geometry}
\usepackage{parskip}
\usepackage{mathtools}
\usepackage{siunitx}
\usepackage{listings}
\usepackage{csquotes}
\usepackage{caption}
\usepackage{subcaption}
\usepackage{comment}
\usepackage{pdfpages}
\usepackage{pict2e}

% \DeclareRobustCommand{\slashcirc}{{\mathpalette\doslashcirc\relax}}

% \makeatletter
% \newcommand\doslashcirc[2]{%
%   \sbox\z@{$#1\m@th\circ$}%
%   \setlength\unitlength{\wd\z@}
%   \begin{picture}(1,1)
%   \roundcap
%   \put(0,0){\box\z@}
%   \put(0,0){\line(1,1){1}}
%   \end{picture}%
% }
% \makeatother


% %% Some packages you will need
% \usepackage{pgfplots}
% \usepackage{pgfplotstable}
% \usepackage{booktabs}
% \usepackage{array}
% \usepackage{colortbl}


\definecolor{arduinoorange}{HTML}{FFA500}
\definecolor{arduinogray}{HTML}{808080}
\definecolor{arduinoblue}{HTML}{007ACC}
\definecolor{arduinogreen}{HTML}{469B00}

\lstset{
  language=C++,
  basicstyle=\ttfamily\footnotesize,
  keywordstyle=\color{arduinoorange},
  stringstyle=\color{arduinogreen},
  commentstyle=\color{arduinogray},
  moredelim=[s][\color{arduinoblue}]{\#}{\ },
  morekeywords={digitalRead,digitalWrite,pinMode,analogRead,analogWrite,Serial,begin,HIGH,LOW},
  frame=tb,
  tabsize=4,
  showstringspaces=false,
  breaklines=true,
  numbers=left,
  numberstyle=\tiny\color{arduinogray},
  numbersep=5pt,
  extendedchars=true,
  literate={á}{{\'a}}1 {ã}{{\~a}}1 {é}{{\'e}}1,
}

\lstdefinestyle{Arduino}
{
  language=C++,
  basicstyle=\ttfamily\footnotesize,
  keywordstyle=\color{arduinoorange},
  stringstyle=\color{arduinogreen},
  commentstyle=\color{arduinogray},
  moredelim=[s][\color{arduinoblue}]{\#}{\ },
  morekeywords={digitalRead,digitalWrite,pinMode,analogRead,analogWrite,Serial,begin},
  frame=tb,
  tabsize=4,
  showstringspaces=false,
  breaklines=true,
  numbers=left,
  numberstyle=\tiny\color{arduinogray},
  numbersep=5pt,
  extendedchars=true,
  literate={á}{{\'a}}1 {ã}{{\~a}}1 {é}{{\'e}}1,
  backgroundcolor=\color{black!85},
  rulecolor=\color{arduinoorange},
  frame=single,
  frameround=tttt,
  framexleftmargin=6pt,
  framexrightmargin=6pt,
  framextopmargin=6pt,
  framexbottommargin=6pt,
  breaklines=true,
  postbreak=\raisebox{0ex}[0ex][0ex]{\ensuremath{\color{red}\hookrightarrow\space}},
}

\usepackage[
    backend=biber,
    backref=true,
    backrefstyle=none,
    sortcites=true,
    sorting=none,
    doi=false, % doi informatie wordt niet weergegeven
    %uniquename=true,
    %uniquelist=true,
    maxcitenames=3,
    %issn=false, werkt niet
    language=american
]{biblatex}
\addbibresource{information/Sources.bib}
\DefineBibliographyStrings{english}{
    backrefpage = {blz.},
    backrefpages = {blz.},
}
\makeglossaries
\definecolor{Grey1}{HTML}{343434}
\graphicspath{{./Media/Figuren/}}
 \geometry{
 a4paper,
 total={170mm,257mm},
 left=20mm,
 top=20mm,
 }
\hypersetup{
    colorlinks=true,
    linkcolor=blue,
    filecolor=magenta,      
    urlcolor=cyan,
    pdftitle={Overleaf Example},
    pdfpagemode=FullScreen,
    }


\begin{document}
\title{
\includegraphics[width=3.5in]{img/Logo/Logo.png} \\
\vspace*{1in}
\textbf{Project 4}\\
\textit{Plan of Approach}\\
Version 1
}
\author{
\vspace*{0.5in} \\
  Written by:\\
  Laurens van der Drift\\
  Justin van der Reijden\\
  Luuk van Kappel\\
  Marnix Harmsen\\
		\vspace*{0.2in} \\
		The Hague University of Applied Sciences\\
        \textbf{Electrical Engineering}\\
        Delft, The Netherlands
       } 
\maketitle
\phantomsection
\section*{Version History} \addcontentsline{toc}{section}{Version History}

\begin{table}[h]
\begin{tabular}{|l|l|l|l|}
\hline
\rowcolor[HTML]{4472C4} 
{\color[HTML]{FFFFFF} \textbf{Version}} &
  {\color[HTML]{FFFFFF} \textbf{Date}} &
  {\color[HTML]{FFFFFF} \textbf{Changes}} &
  {\color[HTML]{FFFFFF} \textbf{Author}} \\ \hline
\rowcolor[HTML]{D9E1F2} 
1.0 &
  \multicolumn{1}{c|}{\cellcolor[HTML]{D9E1F2}05-02-2024} &
 N.v.t. &
  Alset Innovations \\ \hline

% \rowcolor[HTML]{FFFFFF} 
% 2.0 &
%   \multicolumn{1}{c|}{\cellcolor[HTML]{FFFFFF}7-4-2023} &
%  Feedback van FeedbackFruits toegepast &
%   Infra   Vroom \\ \hline

% \rowcolor[HTML]{D9E1F2} 
% 3.0 &
%   \multicolumn{1}{c|}{\cellcolor[HTML]{D9E1F2}4-6-2023} &
%  Bijgewerkt voor Assessment 3 &
%   Infra   Vroom \\ \hline

\end{tabular}
\end{table}
%\addcontentsline{toc}{section}{Verklarende Woordenlijst}
\printglossaries
\newglossaryentry{tender}
{
    name=\textit{tender},
    description={De inschrijving om op een kavel een wind turbine park te bouwen.}
}

\phantomsection
\addcontentsline{toc}{section}{List of Figures}
\listoffigures

\newpage
\tableofcontents
\section{The Assignment}
In the ambitious realm of the \textbf{PROJE4} course, \textbf{Alset Innovations} is tasked with the exciting challenge of designing and implementing a cutting-edge solution for Smart Energy Devices. The focal point of our project revolves around creating a sensored brushless-DC (BLDC) motor controller for an electric bicycle. This endeavor aligns seamlessly with the broader objective of utilizing intelligent embedded technology to enhance efficiency and conserve energy in various applications.

\subsection{Project Scope:}
Our primary goal is to develop a motor controller capable of seamlessly driving a \SI{24}{\volt}, \SI{78}{\watt} BLDC motor in a closed-loop configuration using hall sensor feedback. Our innovative twist lies in the selection of the \textbf{STM32F103C8T6} microcontroller, a new and unexplored avenue, which distinguishes us from our peers who may opt for the more conventional \textbf{RP2040}. This strategic decision places \textbf{Alset Innovations} at the forefront of innovation of this course.

% \subsection{Key Technical Requirements:}
% \begin{enumerate}
%     \item Utilize discrete transistors to control the three motor phases.
%     \item Design a custom printed circuit board (PCB) for the entire system, showcasing our prowess in hardware development.
%     \item Measure and display the current delivered to the motor through the user interface.
%     \item Implement a user-friendly control input (throttle) on the user interface for speed control.
%     \item Integrate the \textbf{STM32F103C8T6} microcontroller to process hall sensor signals, demonstrating our proficiency in microcontroller programming.
%     \item Deliver a comprehensive technical report detailing our design choices, simulations, schematics, and testing results.
% \end{enumerate}

\subsection{Strategic Approach:}
\textbf{Alset Innovations} recognizes the uniqueness of our project by opting for the \textbf{STM32} microcontroller, thereby charting a pioneering course in the implementation of smart energy devices. Our approach involves meticulous planning and execution, encompassing the following key phases:

\begin{enumerate}
    \item \textbf{Initial Research and Exploration:}
        \begin{itemize}
            \item Dive into the specifics of BLDC motor control, \textbf{STM32} microcontroller capabilities, and related electronic circuit theory.
            \item Explore the feasibility of our chosen microcontroller and its potential advantages in comparison to the \textbf{RP2040}.
        \end{itemize}

    \item \textbf{Hardware and PCB Design:}
        \begin{itemize}
            \item Research, compare and test different \textbf{MOSFETS} and \textbf{GATE} drivers.
            \item Develop a detailed schematic and layout for the PCB, focusing on optimal integration and functionality.
            \item Address the challenges posed by components like the \textbf{QFN-56} package of the \textbf{STM32} microcontroller.
        \end{itemize}

    \item \textbf{Software Development:}
        \begin{itemize}
            \item Craft robust code for the microcontroller, ensuring seamless motor control and user interface interaction.
            \item Conduct thorough testing on real prototype hardware, substantiating our software's reliability.
        \end{itemize}

    \item \textbf{Documentation and Reporting:}
        \begin{itemize}
            \item Maintain an exhaustive record of our design process, simulations, and results.
            \item Adhere to the specified template for the final technical report, showcasing our commitment to clarity and precision.
        \end{itemize}
\end{enumerate}

In summary, \textbf{Alset Innovations} stands poised to redefine the landscape of smart energy devices through our unique blend of innovation, strategic decision-making, and meticulous execution. Our Plan of Approach sets the stage for a transformative journey in the world of electrical engineering.

\section{Project Activity}
During this project a number of deadlines must be met four in total. 

Plan of Approach
Test 1
Test 2
Final Test
\section{Project Boundaries}
\section{Costs and revenues}
\section{Organisation}
\section{Planning}
\section{Product}
\section{Quality}
% De lezer leest in dit hoofdstuk hoe je de kwaliteit van de (tussen)producten uit de vorige paragraaf
% waarborgt. Dit kun je bijvoorbeeld beschrijven door:
% -
% - Of en in welke mate je gebruikmaakt van extern advies (bijv. iemand buiten je projectgroep die
% je (tussen)product controleert). Geef aan hoe je de feedback verwerkt;
% - Hoe je omgaat met regelgeving, veiligheid, programma van eisen etc;

% Te beschrijven welke software, materialen en andere hulpbronnen je gebruikt en hoet dit
% bijdraagt aan de kwaliteit van jullie projectresultaat.

% Om de kwaliteit voor het project te waarborgen zal er gebruik gemaakt worden van de software STM32CubeIDE, waarmee de stm32-microcontroller geprogrammeerd kan worden. De stm32-chip is gebruikelijk in de industrie, vanwege zijn manier van programmeren, zoals het instellen van de GPIO en klokfrequenties. Bovendien zou het realtime simulatiepakket Caspoc gebruikt kunnen worden. Hiermee kan de werking en aansturing van de motordriver getest worden. Daarnaast wordt KiCad 7.0 gebruikt voor het ontwerpen van PCB's. Kicad is opensource en beschikt over de functies om een PCB te ontwerpen voor de motordriver. Het bestellen van de PCB wordt bij Eurocircuits gedaan, vanwege afspraken die school heeft gemaakt. \\
% Voor het aanschaffen van materialen zijn bedrijven zoals Conrad, Farnell, Texas Instruments en Mouser geschikt, omdat zij beschikken over een aanbod van hardware voor de motordriver.\\
% Voor het vragen van advies kan er terecht bij:
% \begin{itemize}
%   \item Jesse op den Brouw: C-programmeur
%   \item Stephen O'Loughlin: algemeen advies
%   \item Ad van den Bergh: programmeur
%   \item Diëgo Zuidervliet: kennis voor het aansturen van motordrivers
%   \item Paul Witte: PCB-designer
% \end{itemize}
% Voor het controleren van de (tussen)producten worden assesments gegeven met de docenten die hiervoor een aantal punten geven. Bovendien wordt er elke week een vergadering georganiseerd, waarbij de begeleider betrokken is en eventueel feedback kan geven.\\
% Het ontvangen van feedback wordt meegenomen in de toekomstige activiteiten om zo tot een beter ontwerp te komen. Daarnaast hoeft er geen rekening gehouden te worden met de regelgeving, omdat het eindproduct voor eigengebruik is. Voor de veiligheid moeten componenten gekoeld worden en niet kunnen ontploffen. Het eindproduct zal de minimale eisen hebben volgens het programma van eisen, maar de mogelijkheid bestaat dat er extra functies worden toegevoegd.\\
% Alle bovengenoemde producten moeten zorgen dat er een functionele motordriver ontworpen kan worden.



To ensure the quality of the project, the following software and resources will be utilized:

\begin{itemize}
  \item \textbf{STM32CubeIDE}: This software will be used for programming the STM32 microcontroller, which is commonly used in the industry due to its programming capabilities such as configuring GPIO and clock frequencies.
  \item \textbf{Caspoc Real-time Simulation Package}: This tool will be employed to test the operation and control of the motor driver in real-time simulations.
  \item \textbf{KiCad 7.0}: KiCad, an open-source software, will be utilized for designing PCBs for the motor driver, as it provides the necessary features for PCB design.
  \item \textbf{Eurocircuits}: PCB manufacturing will be done through Eurocircuits as per the agreements made by the school.
\end{itemize}

For procurement of materials, companies such as Conrad, Farnell, Texas Instruments, and Mouser will be considered due to their hardware offerings for the motor driver.

Expert advice can be sought from the following individuals:
\begin{itemize}
  \item Jesse op den Brouw: C programmer
  \item Stephen O'Loughlin: General advice
  \item Ad van den Bergh: Programmer
  \item Diëgo Zuidervliet: Expertise in motor driver control
  \item Paul Witte: PCB designer
\end{itemize}

Assessments will be conducted by instructors to evaluate the (intermediate) products. Additionally, weekly meetings will be held involving the supervisor to provide feedback.

Feedback received will be incorporated into future activities to enhance the design. Regulatory considerations are not necessary as the end product is for personal use. However, components must be cooled for safety, and precautions must be taken to prevent explosions. The final product will meet the minimum requirements outlined in the specifications, with the possibility of additional features being added.

All the resources will contribute to the development of a functional motor driver.


\phantomsection
\addcontentsline{toc}{section}{Referenties}
\printbibliography
\end{document}
